We investigate the spectral mechanism governing mode stabilisation in the
Cosmochrony relaxation framework.
Starting from the admissibility condition defining the projective resolution
$\Lambda_{\mathrm{proj}}$, we derive its dependence on the isoperimetric
capacity of the relaxation graph and show that
$\Lambda_{\mathrm{proj}} \asymp h(G)^2$.
For expander families satisfying spectral-isoperimetric saturation, this
implies $\Lambda_{\mathrm{proj}} \asymp \lambda_2$.

In the Lubotzky--Phillips--Sarnak (LPS) graph model with fixed prime $p$,
the spectral gap converges to a constant, making $\Lambda_{\mathrm{proj}}$
asymptotically static.
In this regime, the stabilisation of modes is governed not by the admissibility
threshold but by saturation of the cumulative spectral count below
$\Lambda_{\mathrm{proj}}$.

Combining this counting mechanism with the representation structure of ADE
substrates yields a factorised prediction for mass ratios,
\[
  \frac{\mathcal{M}_i}{\mathcal{M}_j}
  \;\propto\;
  \frac{F_{\mathrm{KM}}(\lambda_i)}{F_{\mathrm{KM}}(\lambda_j)}
  \cdot
  \frac{\dim\rho_{\lambda_j}}{\dim\rho_{\lambda_i}},
\]
where $F_{\mathrm{KM}}$ is the Kesten--McKay cumulative distribution function.

Numerical evaluation shows that this single-level mechanism produces mass
ratios of order unity and inverts the ordering expected from the admissibility
envelope.
We identify the absence of a dynamical threshold $\Lambda_{\mathrm{proj}}(n)$
as the origin of both limitations and outline directions for restoring
threshold dynamics within the Cosmochrony framework.
